% Owner: TV2 (P5.2a) -- Reviewer: TV3
% Tables T4, T5, and B2B3 are generated by `python -m experiments.report_tables`.
\subsection{Final-Move Selection and UCT Variants}
\label{subsec:results_tv2}

\textbf{Final-move selection on tactical positions (RQ5).} Table~\ref{tab:t4} runs MCTS with a small budget ($c = 0.25$) from three positions whose correct moves are known, $200$ times each with different seeds, and records how often each rule would pick a correct move from the same root statistics. On trap E1, where X must block at $(2,2)$, all three rules are always right: the blocking move collects most of the visits and also has the best mean. On the diagonal fork position, the most visited child is a correct move in $77\%$ of the runs, the child with the best mean in $81\%$, and the two rules disagree in $19\%$. On the $9\times 9$ open-three position, all rules find one of the two blocking moves in only about $13\%$ of the runs: with $80$ iterations and over $40$ candidate moves, the decisive replies are barely sampled, and no selection rule can recover a move that the tree has not explored. The rules differ only when the root statistics are noisy, and there they disagree in up to $30\%$ of the runs.

\begin{table}[t]
\centering
\caption{T4: Percentage of $200$ seeds in which each final-move rule picks a correct move}
\label{tab:t4}
% Generated by experiments/report_tables.py -- do not edit by hand.
\begin{tabular}{lrrrrr}
\toprule
Position & Iter. & Robust & Max & Max-robust & Disagree \\
\midrule
E1 trap ($3\times 3$) & 60 & 100.0 & 100.0 & 100.0 & 0.0 \\
Diagonal fork ($3\times 3$) & 50 & 77.0 & 81.0 & 79.5 & 19.0 \\
Open three ($9\times 9$) & 80 & 12.5 & 13.0 & 12.5 & 30.0 \\
\bottomrule
\end{tabular}

\end{table}

\textbf{Selection rules in play.} Table~\ref{tab:t5} plays MCTS with each rule against perfect Minimax on $3\times 3$ (reporting the share of games not lost, since winning is impossible) and against a fixed MCTS baseline with $60$ iterations on $9\times 9$ (reporting the win rate), with $200$ games per cell. More iterations help every rule, from about $62\%$ to $91$--$96\%$ non-losses on $3\times 3$ and from about $33\%$ to $96.5$--$99.5\%$ wins on $9\times 9$. The differences between the rules are small. Only two of the $27$ comparisons with the max rule are significant at the $5\%$ level, in opposite directions: max-robust beats max on $3\times 3$ with $500$ iterations ($93.5\%$ against $87.5\%$, $p = 0.04$), and max beats robust on $9\times 9$ with $320$ iterations ($99.5\%$ against $96.5\%$, $p = 0.03$). In head-to-head games at $100$ (on $3\times 3$) or $80$ (on $9\times 9$) iterations no rule beats another ($p > 0.2$). Because almost all visits go to the best move as the budget grows (Section~\ref{subsec:math_uct}), the rules converge, and the most visited child is a safe default.

\begin{table}[t]
\centering
\caption{T5: Non-loss rate against perfect Minimax ($3\times 3$) and win rate against MCTS with $60$ iterations ($9\times 9$), in percent, $200$ games per cell}
\label{tab:t5}
\resizebox{\columnwidth}{!}{% Generated by experiments/report_tables.py -- do not edit by hand.
\begin{tabular}{lrrrr}
\toprule
Board & Iter. & Robust & Max & Max-robust \\
\midrule
$3\times 3$ & 50 & 62.0 & 61.5 & 64.5 \\
$3\times 3$ & 100 & 74.5 & 77.5 & 79.5 \\
$3\times 3$ & 250 & 84.0 & 80.5 & 86.5 \\
$3\times 3$ & 500 & 85.5 & 87.5 & 93.5 \\
$3\times 3$ & 1000 & 91.0 & 93.5 & 95.5 \\
\midrule
$9\times 9$ & 40 & 33.0 & 31.0 & 34.0 \\
$9\times 9$ & 80 & 64.5 & 69.5 & 63.5 \\
$9\times 9$ & 160 & 94.5 & 93.0 & 92.5 \\
$9\times 9$ & 320 & 96.5 & 99.5 & 99.5 \\
\bottomrule
\end{tabular}
}
\end{table}

\textbf{UCT variants and improvements (bonus B2, B3).} Table~\ref{tab:b2b3} compares the exploration terms on $9\times 9$ Gomoku against standard UCT with $80$ iterations, with $c = 0.5$ and $200$ games per cell. Removing the square root, or replacing the bonus by $\sqrt{N_p}/(1+N_i)$, weakens MCTS significantly at every budget from $80$ iterations (for example, $72\%$ and $67\%$ against $85\%$ at $160$ iterations): both bonuses decay as $1/N_i$ instead of $1/\sqrt{N_i}$, so a move that loses its first simulations is abandoned too early. Removing the logarithm makes no significant difference at these budgets, which is consistent with the role of the logarithm in Section~\ref{subsec:math_uct}: it matters for the guarantee as $N_p \to \infty$, not for a few hundred iterations. UCB1-Tuned is not better than standard UCT and is worse at $160$ iterations ($77\%$). Progressive bias, which injects the pattern knowledge of the evaluation function into the selection step, is the clear winner: it reaches $65.5\%$ at $80$ iterations, where the standard formula is at $49\%$, and beats standard UCT in $65.5\%$ of head-to-head games ($p < 0.001$). On $3\times 3$ against perfect Minimax, with $1\,000$ iterations the standard formula draws all $200$ games, while removing the logarithm loses $16$ ($p < 0.001$).

\begin{table}[t]
\centering
\caption{B2/B3: Win rate (\%) of UCT variants on $9\times 9$ Gomoku against standard UCT with $80$ iterations; $^{*}$ marks a significant difference from the standard formula at the same budget ($p < 0.05$)}
\label{tab:b2b3}
\resizebox{\columnwidth}{!}{% Generated by experiments/report_tables.py -- do not edit by hand.
\begin{tabular}{lrrrr}
\toprule
Variant & $I=40$ & $I=80$ & $I=160$ & $I=250$ \\
\midrule
Standard & 19.0 & 49.0 & 85.0 & 94.5 \\
No square root & 21.0 & 33.5$^{*}$ & 72.0$^{*}$ & 82.5$^{*}$ \\
No logarithm & 18.5 & 46.0 & 89.0 & 96.5 \\
$\sqrt{N_p}$ bonus & 18.0 & 35.0$^{*}$ & 67.0$^{*}$ & 83.0$^{*}$ \\
UCB1-Tuned & 22.0 & 52.0 & 77.0$^{*}$ & 95.0 \\
Progressive bias & 33.5$^{*}$ & 65.5$^{*}$ & 97.0$^{*}$ & 98.0 \\
\bottomrule
\end{tabular}
}
\end{table}
